\section{Related Work}

% \paratitle{Automated Harness Engineering.} 
% Self-improvement in agentic systems has progressed from optimizing instruction prompts to structured contexts, workflows, and finally the executable harness code that coordinates model reasoning, tool use, memory, and verification \cite{weng2026harness}. Prompt- and context-level methods such as Promptbreeder, GEPA, and ACE, together with workflow-level systems such as ADAS and AFlow, are important precursors, but optimize constrained surfaces of the agent design space \cite{promptbreeder,gepa,ace,adas,aflow}. At the harness-code level, existing methods explore several optimization mechanisms: iterative or self-referential improvement of the scaffold \cite{stop,godelagent,dgm,selfharness}; evolutionary or archive-based search over executable harnesses \cite{metaharness,vero,harnessx}; and trajectory- or observability-guided optimization \cite{ahe,rho}. Joint harness--weight optimization further expands this space by updating the policy model together with its harness \cite{sia}. These methods establish executable harnesses as a high-leverage target for end-to-end optimization. Our benchmark is complementary: rather than proposing another harness optimizer, it fixes the policy model and research protocol and measures how different evolver models carry out long-horizon harness engineering across benchmarks and domains.

\paratitle{Automated Harness Engineering.} 
Recently, agentic self-improvement has evolved from prompt tuning to workflow optimization, and recently, to the automated refinement of the \emph{executable harness code} that orchestrates reasoning, tool use, and memory~\cite{weng2026harness}. While early prompt-level~\cite{promptbreeder,gepa,ace} and workflow-level methods~\cite{adas,aflow} operate on constrained design spaces, recent harness-level engineering leverages diverse mechanisms. These include self-referential scaffold refinement~\cite{stop,godelagent,dgm,selfharness}, evolutionary archive search over executable code~\cite{metaharness,vero}, and trajectory-guided optimization~\cite{harnessx}, alongside recent joint harness--weight co-optimization~\cite{sia}. 
Crucially, rather than proposing a new optimization algorithm, our work provides a complementary orthogonal perspective: we fix the policy model and research protocol to systematically benchmark the intrinsic capacity of frontier foundation models to act as long-horizon harness engineers across heterogeneous domains.

% \paratitle{Agent Benchmarks.}
% Agent benchmarks traditionally measure task execution within a fixed harness \cite{swebench,osworld,taubench,browsecomp,gdpval}, while AI-R\&D benchmarks evaluate agents improving model weights or training pipelines \cite{mlagentbench,mlebench,rebench,paperbench,posttrainbench}. The Meta-Agent Challenge evaluates from-scratch construction of task-specific agent artifacts by their absolute held-out performance \cite{mac}. More recent evaluations move beyond endpoint task scores to study self-evolution itself: SEA-Eval measures longitudinal gain and stability \cite{seaeval}; Harness Updating Is Not Harness Benefit separates the ability to produce updates from the ability to benefit from them \cite{harnessupdating}; EvoAgentBench evaluates cross-task ability transfer \cite{evoagentbench}; and the agent--environment co-evolution survey and SEAGym develop broader measurement frameworks and diagnostics \cite{agentenvcoevolution,seagym}. These works establish how the outcomes and dynamics of self-evolution should be assessed. [NAME] instead benchmarks the foundation model as a long-horizon harness researcher, comparing its ability to improve one shared general-purpose harness across benchmarks and domains under harness-calibrated task construction.


\paratitle{Agent Benchmarks.} 
Existing agent benchmarks predominantly measure task execution within static, fixed harnesses~\cite{swebench,osworld,taubench,browsecomp,gdpval}, whereas AI-R\&D benchmarks evaluate capabilities in optimizing model weights or training pipelines~\cite{mlagentbench,mlebench,rebench,paperbench,posttrainbench}. While the Meta-Agent Challenge evaluates the from-scratch generation of task-specific artifacts~\cite{mac}, an emerging class of evaluations focuses on self-evolutionary dynamics, tracking stability~\cite{seaeval}, distinguishing generation from utilization~\cite{harnessupdating}, assessing capability transfer~\cite{evoagentbench}, or formalizing diagnostic simulators~\cite{seagym}. 
Differing from these frameworks, \textbf{Evo-Bench} shifts the evaluation paradigm. Instead of evaluating task-specific artifacts or unconstrained self-evolution, we explicitly benchmark foundation models as long-horizon research agents tasked with evolving a single, shared general-purpose harness, validated under a rigorous, sensitivity-calibrated task split.